\section{Example: Accelerator Physics of Insertion Devices} \label{ex12}

\W has some feature to investigate aspects of the accelerator physics
of insertion devices. This chapter presents the most important ones.
The results of this calculations really depend on the settings of the
parameters. To get the right answers, check the input and the results
thoroughly.

\subsection{Linear Transfer-Matrix, Multipoles, and Field Maps} \label{ex12a}

Starting point is the Python script \scap{12a}:

\cboxs{\senv \grau \\
python3 ../python/wave\_example\_12a.py\\
cp \$WAVE/input/wave\_example\_12a.in wave.in\\
}\\

The script modifies the file \we \bsp and echos the changes. Please, open
\wi \bsp and study the corresponding variables and their namelists.

Run with this input \W performs three task for an REC-undulator
defined in the namelist \scap{RECN}, see example \ref{ex4}:

\begin{itemize}
\item \underline {Field Map} \\

The variable
\scap{IWBMAP=6}
makes \W writing a field map for later use.
The parameters of the
field map are defined in the namelist \scap{BMAPN}. \W can read and
write field maps of the considered devices. This allows interfacing to
other codes, i.e. importing and exporting magnetic field configurations
between different programs. In the next subsections we will make use
of this field map, mainly to reduce the computing time.

Note: Different namelists might use different units, especially $mm$
and $m$ or even different coordinate systems, respectively.

\begin{figure}[ht]
    \centering
%    \framebox(12.5,8){
      \includegraphics[width=0.65\textwidth]
      {/home/scheer/wdp/wave_example_12a_fig1.pdf}
%    }
    \caption{Transversal distribution of the vertical field in
    the center of the undulator.}
    \label{ex12afig1}
\end{figure}

To visualize the field map, start the plotting GUI and enter in the
Python shell:

\cboxsl{nmap~=~ncread(~''nmap'',~''x:y:z:bx:by:bz'',~''wave\_example12a\_bmap.dat'',~skiphead=6~)\\
nplot(~nmap,~"x:by",~"abs(x)$<$0.028~and~y==0~and~z==0"~)\\
nscan(~nmap,~"x:by",~"abs(x)$<$0.028~and~by$>$1.0~and~y==0~and~z==0"~)\\
nplot(~nmap,~"z*1000:by",~"abs(x+0.014)$<$0.0001"~)\\
setlinecolor('blue')\\
nplot(~nmap,~"z*1000:by",~"abs(x+0.014)$<$0.0001~and~abs(y-0.0052)$<$0.0001",~
"",~"linesame"~)\\
wave\_title()\\
txyz(~"Vert.~Field~at~x~=~0~for~y~=~0/1/2/3/4/5~mm",~"z~[mm]",~~"By~[T]"~)\\
}

The first command creates a Ntuple and reads the field map, skipping
the six header lines. Then the field in the center of the undulator is
plotted. The scan command is used to locate the maximum of the on-axis
field, which is used in the next command. Then the transversal field
distribution is plotted (Fig.~\ref{ex12afig1}) for this corresponding longitudinal position
\scap{x = -0.014 m}, which corresponds to a quarter of the period-length of
the undulator.

Note the dip in the middle, which is a consequence of the slit between
the magnet rows, espacially for \scap{y = 5.2 mm} (blue line). This means that the transveral
field gradients of the device depend on the transversal position of the beam electrons, thus
complicating tracking calculations.

\item \underline{Quadrupole- and Sextupole Terms} \\

\scap{IWBTAB=1} triggers the calculation of the field and the integrals of
the field, the quadrupole- and
sextupole terms etc. along a straight line. For more details refer to the
comments in \wie. In addition the sextupole terms are
written to a \scap{wbtab\_sextupole.dat}\bsp. \scap{IWBTAB=100} does the same, but with respect to the trajectory rather
than a straight line.

By the way: If the field calculation is very
time consuming, in many cases the file written due the
setting of \scap{IWBTAB} is sufficient for later use, e.g. for spectrum
caclulations. In such a case, the field can be determined from the
file via a very fast spline interpolation, see \scap{IRBTAB} in \wi. This
holds of course also for other tabulated fields \scap{B\_y(x)}, e.g. from
measuements. It is also  possible to provide this way two field
components \scap{B\_y(x)} and \scap{Bz(x)}, see \scap{IRBTABZY} in \wi.

\item \underline{Linear Transfer Matrix} \\

\scap{DELTAZ=0.0001} and the other variables of the namelist
\scap{TRALINN} trigger the simple calculation of the linear
transfer-matrix by tracking four principal trajectories through the
device. The results are given in \woe\bsp. The matrix is also written to
\gtis{tralin.wav}\bsp. We copy this file:

\cboxs{cp tralin.wav tralin\_rec.wav}

for the next subsection.

\end{itemize}

The second and third item of this run are not only excecuted for their
results, but also to compare with the later runs, when the field map
will be used. This way, one might get an impression of the accuracy and
uncertaincies related to the use of a field map instead of the
''true'' field.

\subsection{The Linear Transfer-Matrix from the Field Map} \label{ex12b}

We excute the Python script 12b:

\cboxs{\senv \grau \\
python3 ../python/wave\_example\_12b.py\\
cp \$WAVE/input/wave\_example\_12b.in wave.in\\
}\\

and run \W. We compare \gtis{tralin.wav} with the file of the previous
subsection \gtis{tralin\_rec.wav}. Have a look on the matrix terms

\cboxs{m21 = -0.01166\\
m21 = -0.00865}

and

\cboxs{m43 = -0.02564\\
m43 = -0.02500},

which represent the focal strength, i.e. the inverse of the
focal-lenght of the horizontal and vertical focussing. If the required
accuracy is not reached, the grid of the field map must be chosen
finer, or the ''true'' field model must be used.

\subsection{The Kick-Map and Generating Functions} \label{ex12c}

Generating functions (GF) are a poweful tool for tracking calculations,
since they allow symplectic tracking with very few steps through a
whole insertion device. The GF take into account higher order effects,
i.e. they represent the device much better than the linear
transfer-matrix. To calculate a Taylor expansion of the GF $F_2$
numerically, we excute the Python script 12c:

\cboxs{\senv \grau \\
python3 ../python/wave\_example\_12c.py\\
cp \$WAVE/input/wave\_example\_12c.in wave.in\\
}\\

The GF is now calculated in two steps when \W is run. In the first step
the variable \scap{IOPTIC = 1} triggers the subroutine \scap{OPTI}, wchich
tracks a bundle of electrons as defined in the namelist \scap{\$OPTIK}
through the devices. This step can be very time consuming.

In the second step, the GF is calculated. This is triggered by the
variable \scap{IGENFUN = 1} (see also namelist \scap{\$TRANPON}).

Open \wo  and study the output of the
subroutine \scap{TRANPOLY}, which fits the coefficients of the GF.
Please compare the differences between the linear transfer matrices,
quadrupole matrix derived from the GF and matrix calculated in the
subsections above. For real tracking calculations using the GF, one
must check, if the field map is good enough or the ''true'' field
should be used.  And one also should
vary the bunch of tracks, the phase-space etc. and study the
effects on the GF carefully{\ldots} This is beyond the scope of this
example.

%\newpage
We have a look on the tracks used for the GF:

\cboxsl{wplot\\
nlist()\\
ninfo(''n501'')\\
ndump(''n501'',''zi*1000:yi*1000:zpf*1000:ypf*1000'','''',''wave\_example\_12c\_kick-map.dat'')
}\\

The command \gti{nlist()} gives an overview of the Ntuples read from
\scap{WAVE.mhb}.
We note the Ntuple \gti{n501} and print the summary
using the command \gti{ninfo}. Now, we can plot the variables or write
them to a kick-map file using the \gti{ndump} function. The full syntax of
this procedure can be seen by entering:

\cboxs{ndump()}

If no file name is given, the default \gti{ndump.dat} is chosen.

This example also demonstrates a convient way, to extract and convert
data from \scap{WAVE.mhb} for later use. However, since the Ntuples
and histograms are nothing else then \gti{Pandas Dataframes}, you
can directly use the commands provide by the Python package
\gti{Pandas} and \gti{Matplotlib}.

